% =====================================================================
%  Energy-Based Models: Theory and Implementation
%  Drop-in content for the "Classical EBM" and "Langevin Dynamics" sections.
%
%  Requires in the preamble:
%     \usepackage{amsmath, amssymb}
%     \usepackage{listings}      % for the code snippets (optional)
%     \usepackage{booktabs}      % for the mapping tables
%     \usepackage{hyperref}
%
%  If you do not want the code listings, delete the \begin{lstlisting} blocks;
%  the prose and equations stand on their own.
% =====================================================================

\section{Classical Energy-Based Models}

\subsection{Definition: from energy to probability}

An energy-based model (EBM) specifies a probability distribution over a data space
$\mathcal{X} \subseteq \mathbb{R}^{d}$ through a single scalar-valued
\emph{energy function} $E_\theta : \mathcal{X} \to \mathbb{R}$, parameterised by
$\theta$. The associated density follows the Boltzmann--Gibbs form
\begin{equation}
  p_\theta(x) \;=\; \frac{1}{Z(\theta)}\, e^{-E_\theta(x)},
  \qquad
  Z(\theta) \;=\; \int_{\mathcal{X}} e^{-E_\theta(x)}\, dx,
  \label{eq:ebm}
\end{equation}
where $Z(\theta)$ is the \emph{partition function} (normalising constant). The negative
sign in the exponent encodes the defining convention of an EBM: \emph{low energy
corresponds to high probability}. Configurations that resemble the data are assigned low
energy; implausible configurations are assigned high energy. The model is not a
classifier and never emits a label---it emits a single number quantifying how
``data-like'' its input is.

In our implementation the energy $E_\theta$ is a convolutional neural network mapping a
$1\times28\times28$ image to one scalar ($d = 784$). The forward pass \emph{is}
$E_\theta(x)$:
\begin{lstlisting}[language=Python]
# model.py -- CNNModel.forward
def forward(self, x):
    return self.cnn_layers(x).squeeze(dim=-1)  # (B,1,28,28) -> (B,)  one scalar E per image
\end{lstlisting}

\subsection{The intractability of the partition function}

The integral defining $Z(\theta)$ in \eqref{eq:ebm} ranges over the entire
$d$-dimensional space. For $d = 784$ there is no closed form (the integrand is a deep
nonlinear network) and numerical quadrature is infeasible by the curse of dimensionality.
Consequently $Z(\theta)$---and with it the density $p_\theta(x)$ itself---cannot be
evaluated. In particular, the log-likelihood
\begin{equation}
  \log p_\theta(x) \;=\; -\,E_\theta(x) \;-\; \log Z(\theta)
  \label{eq:loglik}
\end{equation}
carries the uncomputable term $\log Z(\theta)$, so na\"ive maximum-likelihood estimation
is unavailable.

Two consequences follow, and they structure the entire method. First, the energy is
identifiable only \emph{up to an additive constant}: shifting $E_\theta \mapsto E_\theta +
c$ leaves $p_\theta$ invariant while rescaling $Z \mapsto Z e^{-c}$. Only \emph{differences}
and \emph{gradients} of the energy are physically meaningful; the absolute level is a
gauge freedom. Second, every subsequent design choice exists to \emph{avoid ever
computing} $Z$. Accordingly, the quantities $Z$, $p_\theta$, and $\log p_\theta$ appear
nowhere in the implementation---their absence is by design.

\subsection{The score function}

The first escape from $Z$ is obtained by differentiating the log-density with respect to
the \emph{input} $x$ (not the parameters). Since $Z(\theta)$ is independent of $x$, its
gradient vanishes:
\begin{equation}
  \nabla_x \log p_\theta(x)
  \;=\; -\,\nabla_x E_\theta(x) \;-\; \underbrace{\nabla_x \log Z(\theta)}_{=\,0}
  \;=\; -\,\nabla_x E_\theta(x).
  \label{eq:score}
\end{equation}
The quantity $s_\theta(x) := \nabla_x \log p_\theta(x) = -\nabla_x E_\theta(x)$ is the
\emph{(Stein) score}. Although the density is inaccessible, its \emph{shape}---the
direction of steepest increase in probability---is a simple backpropagation through the
network, entirely free of $Z$. The score $-\nabla_x E_\theta$ points towards regions of
higher probability (more data-like), and $+\nabla_x E_\theta$ points away from them.

In code, the input gradient $\nabla_x E_\theta$ is provided by the model, with an
implementation-specific stabilising clamp (discussed below):
\begin{lstlisting}[language=Python]
# model.py -- CNNModel.gradient
def gradient(self, x):
    g = super().gradient(x)                     # = grad_x E(x)  via autograd
    if self.grad_clamp is not None:
        g = g.clamp(-self.grad_clamp, self.grad_clamp)
    return g
\end{lstlisting}

It is essential to distinguish the two gradients that arise in EBMs:
$\nabla_x E_\theta$ (with respect to the \emph{input}; free of $Z$; used for sampling) and
$\nabla_\theta \log Z$ (with respect to the \emph{parameters}; the residual difficulty;
addressed in training). The former makes generation tractable; the latter is handled in
Section~\ref{sec:cd}.

% =====================================================================
\section{Langevin Dynamics}
\label{sec:langevin}

\subsection{Sampling via the score}

Knowing the score enables sampling. \emph{Langevin dynamics} is a Markov chain Monte Carlo
scheme whose stationary distribution is exactly $p_\theta$. Starting from an arbitrary
$x_0$, one iterates
\begin{equation}
  x_{t+1} \;=\; x_t \;-\; \frac{\eta}{2}\,\nabla_x E_\theta(x_t)
                        \;+\; \sqrt{\eta}\;\epsilon_t,
  \qquad \epsilon_t \sim \mathcal{N}(0, I),
  \label{eq:langevin}
\end{equation}
which is the Euler--Maruyama discretisation of the overdamped Langevin stochastic
differential equation $dx = -\tfrac12\nabla_x E_\theta(x)\,dt + dW_t$. In the limit
$\eta \to 0$ and $t \to \infty$, the law of $x_t$ converges to $p_\theta$.

The update comprises two competing terms. The \emph{drift} $-\tfrac{\eta}{2}\nabla_x
E_\theta$ is the score of \eqref{eq:score}; it pulls the state downhill in energy, towards
the data manifold (exploitation). The \emph{diffusion} $\sqrt{\eta}\,\epsilon_t$ is an
isotropic Gaussian perturbation; it prevents collapse to a single mode and lets the chain
traverse the distribution (exploration). Absent the diffusion term, the recursion reduces
to deterministic gradient descent and every trajectory converges to the same local
minimum---the failure mode of \emph{mode collapse}. The balance between drift and
diffusion therefore governs sample diversity.

\subsection{Implementation and its deviations from the canonical form}

The sampler realises \eqref{eq:langevin} as a loop; the drift is evaluated through
\texttt{model.gradient}, i.e.\ the clamped score:
\begin{lstlisting}[language=Python]
# sampler.py -- generate_samples (core loop)
for i in range(steps):
    noise_buf.normal_(0, std)        # epsilon ~ N(0, std^2)
    inp_imgs.add_(noise_buf)         #  <-- DIFFUSION (added first; see below)
    inp_imgs.clamp_(min=-1.0, max=1.0)

    grad = model.gradient(inp_imgs)  # grad_x E  (the score, negated)
    inp_imgs.add_(-step_size * grad) #  <-- DRIFT: x <- x - eta * grad_x E
    inp_imgs.clamp_(min=-1.0, max=1.0)
\end{lstlisting}

Our implementation departs from the canonical recursion \eqref{eq:langevin} in three
deliberate respects. These trade mathematical exactness for empirical stability and
diversity, and they should be understood as engineering choices rather than approximations
justified by theory.

\begin{enumerate}
  \item \textbf{Noise precedes the gradient evaluation.} Whereas \eqref{eq:langevin}
    evaluates $\nabla_x E_\theta$ at $x_t$ and then adds noise, the implementation perturbs
    the state first and evaluates the gradient at the perturbed point:
    \begin{equation}
      x_{t+1} \;=\; \bigl(x_t + \sqrt{\eta}\,\epsilon_t\bigr)
                     \;-\; \eta\,\nabla_x E_\theta\bigl(x_t + \sqrt{\eta}\,\epsilon_t\bigr).
      \label{eq:langevin-impl}
    \end{equation}
    This is an alternative, still valid, discretisation of the same SDE. Injecting the
    stochasticity \emph{into} the gradient evaluation decorrelates parallel chains more
    strongly and was found to be decisive in eliminating mode collapse in practice.

  \item \textbf{Decoupled step sizes.} The canonical form couples the drift and diffusion
    scales as $\eta/2$ and $\sqrt{\eta}$, a relationship required for convergence to
    $p_\theta$. The implementation uses two \emph{independent} constants---a large drift
    step \texttt{step\_size} ($\eta = 10$) and a small noise scale \texttt{noise}
    ($\approx 5\times10^{-3}$). The resulting chain is therefore an \emph{approximate}
    sampler, closer in spirit to a stochastically regularised energy minimiser than to an
    asymptotically exact MCMC scheme.

  \item \textbf{Gradient clamping.} The per-step input gradient is clipped elementwise to
    $[-0.03,\,0.03]$ inside \texttt{model.gradient} (see the snippet above). This has no
    counterpart in \eqref{eq:langevin}; it caps the displacement of any single coordinate
    per step, preventing the chain from diverging and---critically---keeping the replay
    buffer (Section~\ref{sec:cd}) free of corrupted samples.
\end{enumerate}

Optionally, the noise scale is \emph{annealed} across the chain, decaying linearly from a
large initial value to a small final one,
\begin{equation}
  \sigma_t \;=\; \sigma_{\mathrm{start}}
    + \frac{t}{T-1}\bigl(\sigma_{\mathrm{end}} - \sigma_{\mathrm{start}}\bigr),
  \label{eq:anneal}
\end{equation}
so that early steps explore broadly (promoting diversity) and late steps settle cleanly
(reducing residual noise). Setting $\sigma_{\mathrm{start}} = \sigma_{\mathrm{end}}$
recovers the constant-noise scheme.

% =====================================================================
\section{Training by Contrastive Divergence}
\label{sec:cd}

\subsection{The maximum-likelihood gradient}

Training seeks parameters maximising the expected log-likelihood of the data,
$\max_\theta \mathbb{E}_{x\sim p_{\mathrm{data}}}[\log p_\theta(x)]$. Differentiating
\eqref{eq:loglik} with respect to $\theta$ retains the partition term, but that term admits
a tractable identity. Using $\nabla_\theta \log Z(\theta) = \frac{1}{Z}\int
e^{-E_\theta}(-\nabla_\theta E_\theta)\,dx = -\,\mathbb{E}_{x\sim p_\theta}[\nabla_\theta
E_\theta(x)]$, one obtains
\begin{equation}
  \nabla_\theta\, \mathbb{E}_{p_{\mathrm{data}}}[\log p_\theta(x)]
  \;=\;
  -\,\mathbb{E}_{x\sim p_{\mathrm{data}}}\!\bigl[\nabla_\theta E_\theta(x)\bigr]
  \;+\; \mathbb{E}_{x\sim p_\theta}\!\bigl[\nabla_\theta E_\theta(x)\bigr].
  \label{eq:mle-grad}
\end{equation}
Equivalently, minimising the negative log-likelihood corresponds to the
\emph{contrastive-divergence} objective
\begin{equation}
  \mathcal{L}_{\mathrm{CD}}
  \;=\;
  \underbrace{\mathbb{E}_{x\sim p_{\mathrm{data}}}\!\bigl[E_\theta(x)\bigr]}_{\text{lower energy of real data}}
  \;-\;
  \underbrace{\mathbb{E}_{x\sim p_\theta}\!\bigl[E_\theta(x)\bigr]}_{\text{raise energy of model samples}}.
  \label{eq:cd}
\end{equation}
The two terms act antagonistically: the positive phase depresses the energy of data
(``positive'') samples, while the negative phase elevates the energy of samples drawn from
the model (``negatives''). At the optimum $p_\theta = p_{\mathrm{data}}$, the two
expectations coincide and the gradient vanishes; the energies of real and generated
batches become statistically indistinguishable. A vanishing training gap is therefore a
\emph{signature of convergence}, not of failure---a point of practical importance when
monitoring training.

The intractable expectation over $p_\theta$ in \eqref{eq:cd} is precisely what Langevin
dynamics supplies. Contrastive divergence approximates it with negatives obtained from a
\emph{finite} number $k$ of Langevin steps rather than from an exact draw:
\begin{equation}
  \mathbb{E}_{x\sim p_\theta}[E_\theta(x)] \;\approx\; E_\theta(x^-),
  \qquad x^- = \mathrm{Langevin}(x_0,\,k).
  \label{eq:cd-approx}
\end{equation}

The loss is realised directly, with the sign convention adopted here (the network returns
$E_\theta$; hence the ordering $E_{\mathrm{real}} - E_{\mathrm{fake}}$):
\begin{lstlisting}[language=Python]
# train.py -- cd_step
fake_imgs = sampler.sample_new_exmps(steps=config["cd_steps"], ...)   # x^- ~ p_theta
real_energy, fake_energy = model(inp_imgs).chunk(2, dim=0)            # E(real), E(fake)
cdiv_loss = real_energy.mean() - fake_energy.mean()                  # L_CD
\end{lstlisting}

\paragraph{Energy-scale regularisation.} Because \eqref{eq:cd} constrains only the
\emph{difference} of energies, the absolute scale---the gauge freedom noted after
\eqref{eq:loglik}---remains unpinned and may drift without bound. A quadratic penalty is
therefore added,
\begin{equation}
  \mathcal{L} \;=\; \mathcal{L}_{\mathrm{CD}}
    \;+\; \alpha\Bigl(\mathbb{E}[E_\theta(x^+)^2] + \mathbb{E}[E_\theta(x^-)^2]\Bigr),
  \label{eq:reg}
\end{equation}
which anchors the energies near zero. This term is a stabiliser and does not arise from
the maximum-likelihood derivation.

\subsection{Persistent negatives: the replay buffer}

A finite $k$ (here $k=60$) is far too small to produce a faithful sample of $p_\theta$ when
each chain is initialised from noise. \emph{Persistent contrastive divergence} circumvents
this by maintaining a buffer of past negatives and resuming the chains from it, so that
across training steps the chains accumulate an effectively unbounded number of Langevin
iterations. To prevent staleness, a small fraction of each batch is re-initialised from
noise:
\begin{equation}
  x_0 \;=\;
  \begin{cases}
    \text{fresh noise}, & \text{with probability } \rho \;(=0.05),\\[2pt]
    \text{resumed from buffer}, & \text{with probability } 1-\rho.
  \end{cases}
  \label{eq:buffer}
\end{equation}
After refinement, the negatives are written back into the buffer. The fidelity of the
entire training signal rests on the buffer remaining populated with high-quality samples;
this is the reason the gradient clamp of Section~\ref{sec:langevin} is indispensable.

\begin{lstlisting}[language=Python]
# sampler.py -- sample_new_exmps
n_new    = np.random.binomial(self.sample_size, self.new_sample_ratio)  # ~5% fresh
rand_imgs = torch.rand((n_new,) + self.img_shape) * 2 - 1               # fresh noise
old_imgs  = torch.cat(random.choices(self.examples, k=self.sample_size - n_new))
inp_imgs  = torch.cat([rand_imgs, old_imgs]).detach().to(self.device)
inp_imgs  = Sampler.generate_samples(self.model, inp_imgs, steps=steps, ...)  # k steps
self.examples = list(inp_imgs.cpu().chunk(self.sample_size)) + self.examples  # persist
self.examples = self.examples[:self.max_len]
\end{lstlisting}

\subsection{Parameter update for a non-stationary objective}

The parameter update is stochastic gradient descent via Adam. Two choices are specific to
the EBM setting. First, the momentum coefficient is set to $\beta_1 = 0$: because the
negatives are drawn from $p_\theta$, the loss surface itself shifts as $\theta$ changes,
rendering the momentum average of \emph{past} gradients stale. With $\beta_1 = 0$, Adam
reduces to an RMSProp-like rule that responds only to the current gradient. Second, the
parameter-gradient norm is clipped to a small value ($0.1$), preventing any single batch
of negatives from inducing a destabilising update:
\begin{equation}
  g \;\leftarrow\; g\,\min\!\Bigl(1,\, \tfrac{c}{\|g\|}\Bigr),\quad c = 0.1,
  \qquad
  \theta \;\leftarrow\; \mathrm{Adam}_{\beta_1=0}(\theta,\, g).
  \label{eq:update}
\end{equation}
We emphasise that this parameter clip (on $\nabla_\theta \mathcal{L}$) is distinct from the
input-gradient clamp of Section~\ref{sec:langevin} (on $\nabla_x E_\theta$): they act on
different gradients, in different places, and serve different purposes.

\begin{lstlisting}[language=Python]
# train.py -- main loop
optimizer.zero_grad()
loss.backward()                                                    # grad_theta L
torch.nn.utils.clip_grad_norm_(model.parameters(), config["grad_clip"])  # ||.|| <= 0.1
optimizer.step()                                                   # theta <- Adam update
scheduler.step()                                                   # lr *= 0.97 per epoch
\end{lstlisting}

% =====================================================================
\section{The Complete Algorithm}

The training procedure interleaves sampling and learning. Each step draws negatives from
the persistent buffer \eqref{eq:buffer}, refines them by $k$ Langevin iterations
\eqref{eq:langevin-impl} using the clamped score, evaluates the regularised
contrastive-divergence loss \eqref{eq:cd}--\eqref{eq:reg}, and updates $\theta$ under the
non-stationary-aware rule \eqref{eq:update}. Crucially, the Langevin chain runs \emph{inside}
each parameter update: sampling is nested within training. Sampling and learning are the
same machinery directed to two ends---the input gradient shapes the samples, the parameter
gradient shapes the landscape---bound together by the single energy function $E_\theta$.

\begin{table}[h]
\centering
\begin{tabular}{@{}lll@{}}
\toprule
Stage & Equation & Location \\
\midrule
Energy $E_\theta$            & \eqref{eq:ebm}                 & \texttt{model.py:forward} \\
Score $-\nabla_x E$          & \eqref{eq:score}               & \texttt{model.py:gradient} \\
Langevin refinement          & \eqref{eq:langevin-impl}       & \texttt{sampler.py:generate\_samples} \\
Replay buffer                & \eqref{eq:buffer}              & \texttt{sampler.py:sample\_new\_exmps} \\
CD loss $+$ regulariser      & \eqref{eq:cd},\,\eqref{eq:reg} & \texttt{train.py:cd\_step} \\
Parameter update             & \eqref{eq:update}              & \texttt{train.py:main} \\
Partition function $Z$       & ---                            & (never computed, by design) \\
\bottomrule
\end{tabular}
\caption{Correspondence between the mathematical formulation and the implementation.}
\end{table}

% =====================================================================
\section{Inference (Generation)}
\label{sec:inference}

Generation is Langevin dynamics on the trained, \emph{frozen} energy. Starting from noise,
one iterates \eqref{eq:langevin-impl} for $T$ steps with $\theta$ fixed:
\begin{equation}
  x_0 \sim \mathcal{U}(-1,1)^d,
  \qquad
  x_{t+1} = \mathrm{Langevin}(x_t),\quad t = 0,\dots,T-1,
  \qquad \hat{x} = x_T.
  \label{eq:inference}
\end{equation}
No new mechanism is introduced: inference invokes the \emph{same} sampler used to generate
training negatives, with parameter gradients disabled. It therefore exercises only the
sampling half of the method (the score and Langevin dynamics), not the learning half.

\begin{lstlisting}[language=Python]
# inference.py
start   = torch.rand((N_SAMPLES,) + img_shape) * 2 - 1     # x_0 ~ noise
samples = Sampler.generate_samples(model, start, steps=N_STEPS,
                                   step_size=STEP_SIZE, noise=NOISE, ...)
\end{lstlisting}

Three conditions distinguish inference from the sampling performed during training, and
they account for the empirically observed difficulty of generation. (i) \emph{Cold start:}
training negatives resume from a buffer carrying the history of many prior steps, whereas
inference begins from fresh noise with only $T$ steps available; the energy landscape is
navigable from the buffer's occupied regions but not necessarily from arbitrary noise.
(ii) \emph{Frozen parameters:} during inference the landscape is static and only the input
evolves. (iii) \emph{Free hyperparameters:} the step size, noise scale, and step count may
be chosen at will; an excessive ratio of drift to diffusion collapses the chain to a single
mode, formally when $\eta\|\nabla_x E\| \gg \sqrt{\eta}\,\|\epsilon\|$. Sample diversity may
be quantified by the mean pairwise distance $\frac{1}{N^2}\sum_{i,j}\|x_i - x_j\|_2$, which
approaches zero under collapse.

In summary, Sections~\ref{sec:langevin}--\ref{sec:cd} construct and train the energy
$E_\theta$, while inference (Section~\ref{sec:inference}) applies the sampling subroutine
alone to the frozen result. The same energy function that assigns scores during training is
the one descended by Langevin dynamics at generation time.
